(1) 取 \(C\) 为正实数，利用三维球坐标体积元 \(d^3r=4\pi r^2dr\)，归一化条件是
\[
1=4\pi |C|^2\int_0^\infty r^2e^{-2r/a_Z}dr
=4\pi|C|^2\frac{2!}{(2/a_Z)^3}
=\pi a_Z^3|C|^2.
\]
因此 \(C=1/\sqrt{\pi a_Z^3}\)。径向概率密度为
\[
P_r(r)=4\pi r^2|\psi_{1s}(r)|^2
=\frac{4r^2}{a_Z^3}e^{-2r/a_Z},\qquad r>0.
\]
将其对 \(r\) 求导，非零极值满足 \(2/r-2/a_Z=0\)，故最概然半径为 \(r_{\rm mode}=a_Z\)。

(2) 对任意 \(n>-3\)，利用 Gamma 积分
\[
\langle r^n\rangle=
\frac4{a_Z^3}\int_0^\infty r^{n+2}e^{-2r/a_Z}dr
=\frac{(n+2)!}{2^{n+1}}a_Z^n
\]
（对非整数 \(n\) 把阶乘理解为 Gamma 函数）。于是
\[
\langle r\rangle=\frac{3a_Z}{2},\quad
\langle r^2\rangle=3a_Z^2,\quad
\left\langle\frac1r\right\rangle=\frac1{a_Z}.
\]
注意最概然半径 \(a_Z\) 不等于平均半径 \(3a_Z/2\)。

(3) 基态束缚能为
\[
E_{1s}=-\frac{\hbar^2}{2\mu a_Z^2}
=-\frac{\mu Z^2e^4}{2(4\pi\epsilon_0)^2\hbar^2}.
\]
库仑势能的期望值由上一步得到
\[
\langle U\rangle=-\frac{Ze^2}{4\pi\epsilon_0}
\left\langle\frac1r\right\rangle
=-\frac{\hbar^2}{\mu a_Z^2}=2E_{1s}.
\]
再由 \(E=\langle K\rangle+\langle U\rangle\) 得
\[
\langle K\rangle=-E_{1s}=\frac{\hbar^2}{2\mu a_Z^2}>0.
\]
从而 \(2\langle K\rangle+\langle U\rangle=0\)，与 \(U\propto-1/r\) 的量子维里定理完全一致。

(4) 球内累计概率为
\[
\Pr(r<R)=\int_0^R\frac{4r^2}{a_Z^3}e^{-2r/a_Z}dr
=1-e^{-2R/a_Z}\left(1+\frac{2R}{a_Z}+\frac{2R^2}{a_Z^2}\right).
\]
当 \(R\to0\) 该概率趋零、当 \(R\to\infty\) 趋一。固定约化质量与基本常数时 \(a_Z\propto Z^{-1}\)，平均半径和最概然半径都随 \(Z\) 增大而缩小，而束缚能的绝对值 \(|E_{1s}|\propto Z^2\) 增长；若考虑超高核电荷的相对论效应，则非相对论模型需修正。
